\createtitlepage{Homework 9}{June 5, 2026}

\begin{problem}[14.4]
  Theorem 14.24 states that every p-group has nontrivial center. Find the center of $D_8$.
\end{problem}

\begin{solution}
  Recall that $D_8$ has order 16 and is generated by a rotation $\rho$ of order 8 and a reflection $\mu$ with $\mu\rho\mu^{-1} = \rho^{-1}$. An element is in $Z(D_8)$ if and only if it commutes with both $\rho$ and $\mu$.

  For powers of $\rho$: every $\rho^k$ commutes with $\rho$. We check which commute with $\mu$: $\mu\rho^k = \rho^{-k}\mu$, so $\rho^k\mu = \mu\rho^k$ if and only if $\rho^k = \rho^{-k}$, i.e., $\rho^{2k} = e$, i.e., $8 \mid 2k$, i.e., $4 \mid k$. So among rotations, only $\rho^0 = e$ and $\rho^4$ are central.

  For elements of the form $\mu\rho^k$: we check $(\mu\rho^k)\rho = \mu\rho^{k+1}$ and $\rho(\mu\rho^k) = \rho\mu\rho^k = \mu\rho^{-1}\rho^k = \mu\rho^{k-1}$. These are equal only if $\rho^{k+1} = \rho^{k-1}$, i.e., $\rho^2 = e$, which is false. So no reflection is central.

  Therefore $Z(D_8) = \{e, \rho^4\} \cong \Z_2$.
\end{solution}

\begin{problem}[14.5]
  Find the center of $D_7$.
\end{problem}

\begin{solution}
  The group $D_7$ has order 14 and is generated by a rotation $\rho$ of order 7 and a reflection $\mu$ with $\mu\rho\mu^{-1} = \rho^{-1}$.

  For rotations $\rho^k$: $\mu\rho^k = \rho^{-k}\mu$, so $\rho^k$ commutes with $\mu$ if and only if $\rho^k = \rho^{-k}$, i.e., $\rho^{2k} = e$, i.e., $7 \mid 2k$. Since $\gcd(2, 7) = 1$, this requires $7 \mid k$, so only $k = 0$, i.e., only $e$.

  For reflections $\mu\rho^k$: as in the previous problem, $(\mu\rho^k)\rho = \mu\rho^{k+1}$ while $\rho(\mu\rho^k) = \mu\rho^{k-1}$, which are unequal since $\rho \neq e$. So no reflection is central.

  Therefore $Z(D_7) = \{e\}$, which is consistent with $D_7$ not being a $p$-group (since $|D_7| = 14$ is not a prime power).
\end{solution}

\begin{problem}[14.6]
  Let $G = X = S_3$ and make $X$ a $G$-set using conjugation. That is, $*(\sigma, \tau) = \sigma\tau\sigma^{-1}$. Find all the orbits of $X$ using this action. (Write permutations in disjoint cycle notation.)
\end{problem}

\begin{solution}
  The orbits under conjugation are precisely the conjugacy classes of $S_3$. Elements of $S_3$ in disjoint cycle notation are:
  \[%
    S_3 = \{\iota,\, (1\;2\;3),\, (1\;3\;2),\, (1\;2),\, (1\;3),\, (2\;3)\}
  .\]%
  Conjugate permutations have the same cycle type, so the conjugacy classes are determined by cycle type:
  \begin{itemize}
    \item Identity: $\{\iota\}$.
    \item 3-cycles: $\{(1\;2\;3),\, (1\;3\;2)\}$.
    \item Transpositions: $\{(1\;2),\, (1\;3),\, (2\;3)\}$.
  \end{itemize}
  These three sets are the orbits of $X$ under conjugation by $G$.
\end{solution}

\begin{problem}[14.7]
  Let $G = D_4$ and $X$ be the set of all subgroups of $D_4$ with order two. The set $X$ is a $G$-set using conjugation, $*(\sigma, H) = \sigma H\sigma^{-1}$. Find all the orbits of this group action.
\end{problem}

\begin{solution}
  Using standard notation, $D_4 = \{\iota, \rho, \rho^2, \rho^3, \mu, \mu\rho, \mu\rho^2, \mu\rho^3\}$ where $\rho$ is rotation by $\pi/2$ and $\mu$ is a reflection. The subgroups of order two are exactly the subgroups generated by elements of order two:
  \[%
    X = \{\langle \rho^2 \rangle,\, \langle \mu \rangle,\, \langle \mu\rho^2 \rangle,\, \langle \mu\rho \rangle,\, \langle \mu\rho^3 \rangle\}
  .\]%
  The subgroup $\langle \rho^2 \rangle = \{e, \rho^2\}$ is normal in $D_4$ since $\rho^2$ is central, so its conjugacy class is $\{\langle \rho^2 \rangle\}$.

  For the reflections: $\mu$ and $\mu\rho^2$ are reflections across the diagonals, while $\mu\rho$ and $\mu\rho^3$ are reflections across the lines joining midpoints of opposite sides. Conjugating by $\rho$ sends $\mu \mapsto \rho\mu\rho^{-1} = \mu\rho^2$, and $\mu\rho \mapsto \rho(\mu\rho)\rho^{-1} = \mu\rho^3$, so:
  \begin{itemize}
    \item $\{\langle \mu \rangle,\, \langle \mu\rho^2 \rangle\}$ is one orbit.
    \item $\{\langle \mu\rho \rangle,\, \langle \mu\rho^3 \rangle\}$ is one orbit.
  \end{itemize}
  Therefore the orbits are $\{\langle \rho^2 \rangle\}$, $\{\langle \mu \rangle, \langle \mu\rho^2 \rangle\}$, and $\{\langle \mu\rho \rangle, \langle \mu\rho^3 \rangle\}$.
\end{solution}

\begin{problem}[14.8]
  Let $G = U = \{z \in \C | |z| = 1\}$ be the circle group. Then $X = \C$, the set of complex numbers, is a $G$-set with group action given by complex number multiplication. That is, if $z \in U$ and $w \in \C$, $*(z, w) = zw$. Find all the orbits of this action. Also, find $X_G$.
\end{problem}

\begin{solution}
  For $w \in \C$, the orbit of $w$ is
  \[%
    Gw = \{zw \mid z \in U\} = \{zw \mid |z| = 1\}
  .\]%
  Since $|zw| = |z||w| = |w|$, every element of $Gw$ has modulus $|w|$. Conversely, if $w' \in \C$ with $|w'| = |w|$, then $z = w'/w \in U$ (when $w \neq 0$) satisfies $zw = w'$. Thus:
  \begin{itemize}
    \item If $w \neq 0$, the orbit of $w$ is the circle $\{w' \in \C \mid |w'| = |w|\}$ of radius $|w|$ centered at the origin.
    \item If $w = 0$, the orbit is $\{0\}$.
  \end{itemize}
  So the orbits are $\{0\}$ and the circles centered at the origin of each positive radius $r > 0$.

  For $X_G$, an element $w \in \C$ satisfies $zw = w$ for all $z \in U$ if and only if $(z-1)w = 0$ for all $z \in U$, which (taking $z \neq 1$) forces $w = 0$. Hence $X_G = \{0\}$.
\end{solution}

\begin{problem}[14.10]
  Let $G$ be a group of order 9 and suppose that $|X| = 10$. For each possible action of $G$ on $X$, give a list of the orbit sizes. List the orbit sizes from largest to smallest.
\end{problem}

\begin{solution}
  By Theorem 14.17, each orbit size divides $|G| = 9$, so orbit sizes lie in $\{1, 3, 9\}$, and they must sum to $|X| = 10$. The possible partitions of 10 into parts from $\{1, 3, 9\}$ are:
  \begin{itemize}
    \item $9, 1$
    \item $3, 3, 3, 1$
    \item $3, 3, 1, 1, 1, 1$
    \item $3, 1, 1, 1, 1, 1, 1, 1$
    \item $1, 1, 1, 1, 1, 1, 1, 1, 1, 1$
  \end{itemize}
  Each of these is realizable, and these are all possible orbit size lists.
\end{solution}

\begin{problem}[14.19]
  Let $X$ be a $G$-set. Show that $G$ acts faithfully on $X$ if and only if no two distinct elements of $G$ have the same action on each element of $X$.
\end{problem}

\begin{solution}
  Recall that $G$ acts faithfully on $X$ means the kernel of the homomorphism $\phi : G \to S_X$ defined by $\phi(g) = \sigma_g$ is trivial, i.e., the only element of $G$ fixing every $x \in X$ is $e$.

  Suppose $G$ acts faithfully and let $g_1, g_2 \in G$ have the same action on every $x \in X$, i.e., $g_1 x = g_2 x$ for all $x \in X$. Then $g_2^{-1}g_1 x = x$ for all $x \in X$, so $g_2^{-1}g_1$ fixes every element of $X$. By faithfulness, $g_2^{-1}g_1 = e$, hence $g_1 = g_2$.

  Conversely, suppose no two distinct elements of $G$ have the same action on every $x \in X$. Let $g \in G$ fix every $x \in X$, so $gx = ex$ for all $x \in X$. By hypothesis, $g = e$. Hence the kernel of $\phi$ is trivial, so $G$ acts faithfully.
\end{solution}

\begin{problem}[14.22]
  Let $\{X_i | i \in I\}$ be a disjoint collection of sets, so $x_i \cap X_j = \emptyset$ for $i \ne j$. Let each $X_i$ be a $G$-set for the same group $G$.
  \begin{enumerate}
    \item Show that $\bigcap_{i\in I} X_i$ can be viewed in a natural way as a $G$-set, the union of the $G$-sets $X_i$.
    \item Show that every $G$-set $X$ is the union of its orbits.
  \end{enumerate}
\end{problem}

\begin{solution}[(i)]
  Let $X = \bigcup_{i \in I} X_i$. Since the $X_i$ are pairwise disjoint, every $x \in X$ belongs to a unique $X_i$. Define the action of $G$ on $X$ by $gx = g \cdot_i x$, where $\cdot_i$ denotes the given action of $G$ on $X_i$ and $i$ is the unique index with $x \in X_i$. We verify the two conditions of Definition 14.3:
  \begin{enumerate}
    \item For any $x \in X_i$, $ex = e \cdot_i x = x$ since $X_i$ is a $G$-set.
    \item For $g_1, g_2 \in G$ and $x \in X_i$, we have $g_2 x = g_2 \cdot_i x \in X_i$ (since $X_i$ is a $G$-set), and then $(g_1g_2)x = (g_1g_2) \cdot_i x = g_1 \cdot_i (g_2 \cdot_i x) = g_1(g_2 x)$.
  \end{enumerate}
  Hence $X$ is a $G$-set.
\end{solution}

\begin{solution}[(ii)]
  Let $X$ be a $G$-set. By Theorem 14.15, the relation $x_1 \sim x_2$ if and only if $gx_1 = x_2$ for some $g \in G$ is an equivalence relation on $X$, whose equivalence classes are the orbits. Since equivalence classes partition $X$, every element of $X$ belongs to exactly one orbit, and $X$ is the union of its orbits.
\end{solution}

\begin{problem}[14.29]
  Prove that if $G$ is a group of order $p^3$, where $p$ is a prime number, then $|Z(G)|$ is either $p$ or $p^3$. Give an example where $|Z(G)| = p$ and an example where $|Z(G)| = p^3$.
\end{problem}

\begin{solution}
  Since $|G| = p^3$, $G$ is a $p$-group, so by Theorem 14.24 the center $Z(G)$ is nontrivial. By Lagrange's theorem, $|Z(G)| \in \{p, p^2, p^3\}$. Suppose for contradiction that $|Z(G)| = p^2$. Then $|G/Z(G)| = p$, so $G/Z(G)$ is cyclic. But if $G/Z(G)$ is cyclic, then $G$ is abelian, forcing $Z(G) = G$ and $|Z(G)| = p^3$, a contradiction. Hence $|Z(G)| \neq p^2$, and so $|Z(G)|$ is either $p$ or $p^3$.

  Now, for the examples. The abelian group $G = \Z_{p^3}$ satisfies $Z(G) = G$, so $|Z(G)| = p^3$. For $p = 2$, the dihedral group $D_4$ has order $8 = 2^3$ and center $Z(D_4) = \{e, \rho^2\} \cong \Z_2$, so $|Z(D_4)| = 2 = p$. For odd $p$, the Heisenberg group modulo $p$,
  \[%
    G = \left\{ \begin{pmatrix} 1 & a & b \\ 0 & 1 & c \\ 0 & 0 & 1 \end{pmatrix} \;\middle|\; a, b, c \in \Z_p \right\}
  ,\]%
  has order $p^3$ and center $Z(G) \cong \Z_p$, giving $|Z(G)| = p$.
\end{solution}

\begin{problem}[16.2]
  Let $\phi : \Z_{18} \to \Z_{12}$ be the homomorphism where $\phi(1) = 10$.
  \begin{enumerate}
    \item Find the kernel $K$ of $\phi$.
    \item List the cosets in $\Z_{18}/K$, showing the elements in each coset.
    \item Find the group $\phi[\Z_{18}]$.
    \item Give the correspondence between $\Z_{18}/K$ and $\phi[\Z_{18}]$ given by the map $\mu$ described in Theorem 16.2.
  \end{enumerate}
\end{problem}

\begin{solution}[(i)]
  Since $\phi$ is a homomorphism with $\phi(1) = 10$, we have $\phi(n) = 10n \pmod{12}$. The kernel is
  \[%
    K = \{n \in \Z_{18} \mid 10n \equiv 0 \pmod{12}\}.
  \]%
  We need $12 \mid 10n$, i.e., $6 \mid 5n$. Since $\gcd(5,6) = 1$, this requires $6 \mid n$. The elements of $\Z_{18}$ divisible by 6 are $\{0, 6, 12\}$, so $K = \{0, 6, 12\}$.
\end{solution}

\begin{solution}[(ii)]
  Since $|K| = 3$, the quotient $\Z_{18}/K$ has $18/3 = 6$ cosets:
  \begin{align*}
    0 + K &= \{0, 6, 12\}, \\
    1 + K &= \{1, 7, 13\}, \\
    2 + K &= \{2, 8, 14\}, \\
    3 + K &= \{3, 9, 15\}, \\
    4 + K &= \{4, 10, 16\}, \\
    5 + K &= \{5, 11, 17\}
  .\qedhere\end{align*}
\end{solution}

\begin{solution}[(iii)]
  The image is $\phi[\Z_{18}] = \{10n \pmod{12} \mid n \in \Z_{18}\}$. Since $\gcd(10, 12) = 2$, the image is the unique subgroup of $\Z_{12}$ of order $12/\gcd(10,12) = 6$, namely $\phi[\Z_{18}] = \langle 2 \rangle = \{0, 2, 4, 6, 8, 10\} \le \Z_{12}$.
\end{solution}

\begin{solution}[(iv)]
  By the First Isomorphism Theorem, the isomorphism $\mu : \Z_{18}/K \to \phi[\Z_{18}]$ is given by $\mu(n + K) = \phi(n) = 10n \pmod{12}$:
  \[%
    \mu(0 + K) = 0, \quad \mu(1 + K) = 10, \quad \mu(2 + K) = 8, \quad \mu(3 + K) = 6, \quad \mu(4 + K) = 4, \quad \mu(5 + K) = 2
  .\qedhere\]%
\end{solution}
